\chapter{Introduction}

Search algorithms are fundamental components of data manipulation libraries, providing efficient methods to locate specific elements within datasets. The efficiency of these algorithms becomes increasingly important as dataset sizes grow, affecting the overall performance of applications that rely on frequent search operations.

The Dataclenz library implements several search algorithms, with this study focusing on two prominent options: Binary Search and Jump Search. Both are designed for sorted datasets but differ significantly in their approach and theoretical time complexity.

\section{Background}

Search algorithms can be broadly categorized based on their time complexity and requirements. Linear Search (O(n)) examines each element sequentially, making it suitable for unsorted data but inefficient for large datasets. In contrast, algorithms like Binary Search (O(log n)) and Jump Search (O(sqrt(n))) require sorted data but offer significant performance improvements for large datasets.

\section{Motivation}

While theoretical time complexities provide general performance expectations, real-world factors such as data types, memory access patterns, and implementation details can influence actual performance. This research aims to bridge the gap between theoretical understanding and practical application by conducting empirical tests on varied datasets.

\section{Research Questions}

This study addresses the following key questions:

\begin{enumerate}
    \item How do Binary Search and Jump Search performance metrics compare across different dataset sizes?
    \item What impact do different data types (integers, floats, strings) have on search algorithm performance?
    \item How do memory usage patterns differ between these algorithms?
    \item At what dataset sizes does the performance difference between algorithms become significant?
    \item How closely do empirical results align with theoretical time complexity expectations?
\end{enumerate}

\section{Contribution}

This research contributes to the field of algorithm analysis by:

\begin{itemize}
    \item Providing empirical evidence of algorithm performance across varied conditions
    \item Identifying practical thresholds for algorithm selection
    \item Demonstrating the impact of data types on search performance
    \item Offering insights into the relationship between theoretical complexity and real-world behavior
\end{itemize}

The findings from this study will help developers make informed decisions when selecting search algorithms for specific applications, considering factors beyond theoretical time complexity alone.